% Part: first-order-logic
% Chapter: natural-deduction
% Section: proving-things-quant

\documentclass[../../../include/open-logic-section]{subfiles}

\begin{document}

\olfileid[hi]{fol}{ntd}{prq}

\olsection{परिमाणकों वाली \usetoken{p}{derivation}}

\begin{ex}
परिमाणकों के साथ काम करते समय यह सुनिश्चित करना होता है कि अभिलाक्षणिक-चर
शर्त का उल्लंघन न हो। इसके लिए कभी-कभी कुछ अनुमानों को लगाने का क्रम बदलकर
देखना पड़ता है। सामान्यतः अभिलाक्षणिक-चर शर्त के अधीन नियमों को पहले सँभालना
सहायक होता है (पूरे प्रमाण में वे नीचे आएँगे)।

आइए देखें कि !!a{derivation} कैसे दी जा सकती है—इस !!{formula} की:
$\lexists[x][\lnot !A(x)] \lif \lnot \lforall[x][!A(x)]$.
सामान्य रीति से आरम्भ करते हुए लिखते हैं:
\begin{prooftree}
\AxiomC{}
\UnaryInfC{$\lexists[x][\lnot !A(x)]\lif \lnot \lforall[x][!A(x)]$}
\end{prooftree}
सबसे पहले लिखते हैं कि उस अंतिम चरण को \Intro{\lif} नियम से उचित ठहराने के
लिए क्या चाहिए।
\begin{prooftree}
\AxiomC{$\Discharge{\lexists[x][\lnot !A(x)]}{1}$}
\DeduceC{$\lnot \lforall[x][!A(x)]$}
\DischargeRule{\Intro{\lif}}{1}
\UnaryInfC{$\lexists[x][\lnot !A(x)]\lif \lnot \lforall[x][!A(x)]$}
\end{prooftree}
चूँकि $\lnot \lforall[x][!A(x)]$ पर लगाने के लिए कोई स्पष्ट नियम नहीं है,
इसलिए !!{derivation} को इस प्रकार व्यवस्थित करेंगे कि \Elim{\lexists} नियम
लगा सकें। यहाँ अभिलाक्षणिक-चर शर्त पर ध्यान देकर ऐसा अचर चुनना होगा जो
$\lexists[x][!A(x)]$ में या उन मान्यताओं में न आए जिन पर वह निर्भर है।
(फिर भी, चूँकि कोई !!{constant} नहीं आया है, कोई भी चुनाव ठीक रहेगा।)
\begin{prooftree}
\AxiomC{$\Discharge{\lexists[x][\lnot !A(x)]}{1}$}
\AxiomC{$\Discharge{\lnot !A(a)}{2}$}
\DeduceC{$\lnot \lforall[x][!A(x)]$}
\DischargeRule{\Elim{\lexists}}{2}
\BinaryInfC{$\lnot \lforall[x][!A(x)]$}
\DischargeRule{\Intro{\lif}}{1}
\UnaryInfC{$\lexists[x][\lnot !A(x)] \lif \lnot \lforall[x][!A(x)]$}
\end{prooftree}
$\lnot \lforall[x][!A(x)]$ को व्युत्पन्न करने के लिए \Intro{\lnot} नियम
लगाने का प्रयत्न करेंगे। इसके लिए कोई अंतर्विरोध व्युत्पन्न करना होगा, जिसमें
$\lforall[x][!A(x)]$ को अतिरिक्त मान्यता के रूप में लिया जा सकता है। निस्संदेह
इस अंतर्विरोध में मान्यता $\lnot !A(a)$ आ सकती है, जिसे \Elim{\lexists} अनुमान
मुक्त करेगा। इसे इस प्रकार व्यवस्थित कर सकते हैं:
\begin{prooftree}
\AxiomC{$\Discharge{\lexists[x][\lnot !A(x)]}{1}$}
\AxiomC{$\Discharge{\lnot !A(a)}{2}, \Discharge{\lforall[x][!A(x)]}{3}$}
\DeduceC{$\lfalse$}
\DischargeRule{\Intro{\lnot}}{3}
\UnaryInfC{$\lnot \lforall[x][!A(x)]$}
\DischargeRule{\Elim{\lexists}}{2}
\BinaryInfC{$\lnot \lforall[x][!A(x)]$}
\DischargeRule{\Intro{\lif}}{1}
\UnaryInfC{$\lexists[x][\lnot !A(x)]\lif \lnot \lforall[x][!A(x)]$}
\end{prooftree}
लगता है कि हम अंतर्विरोध पाने के निकट हैं। लगाने के लिए सबसे सरल नियम
\Elim{\lforall} है, जिस पर अभिलाक्षणिक-चर शर्त लागू नहीं होती। सार्वपरिमाणित
चर~$x$ के स्थान पर कोई भी इच्छित पद रख सकते हैं; अतः अंतर्विरोध तक पहुँचने के
लिए~$a$ का ही प्रयोग जारी रखना सबसे उपयुक्त है।
\begin{prooftree}
\AxiomC{$\Discharge{\lexists[x][\lnot !A(x)]}{1}$}
\AxiomC{$\Discharge{\lnot !A(a)}{2}$}
\AxiomC{$\Discharge{\lforall[x][!A(x)]}{3}$}
\RightLabel{\Elim{\lforall}}
\UnaryInfC{$!A(a)$}
\RightLabel{\Elim{\lnot}}
\BinaryInfC{$\lfalse$}
\DischargeRule{\Intro{\lnot}}{3}
\UnaryInfC{$\lnot \lforall[x][!A(x)]$}
\DischargeRule{\Elim{\lexists}}{2}
\BinaryInfC{$\lnot \lforall[x][!A(x)]$}
\DischargeRule{\Intro{\lif}}{1}
\UnaryInfC{$\lexists[x][\lnot !A(x)]\lif \lnot \lforall[x][!A(x)]$}
\end{prooftree}

विशेषतः परिमाणकों के साथ काम करते समय इस बिंदु पर दोबारा जाँचना आवश्यक है
कि अभिलाक्षणिक-चर शर्त का उल्लंघन तो नहीं हुआ। हमने इस शर्त के अधीन केवल
\Elim{\exists} नियम लगाया है, और अभिलाक्षणिक चर~$a$ उन किसी भी मान्यताओं में
नहीं आता जिन पर वह निर्भर है; अतः यह सही !!{derivation} है।
\end{ex}

\begin{ex}
कभी-कभी हम !!a{formula} को अन्य !!{formula}s से व्युत्पन्न करते हैं। ऐसी दशाओं
में अमुक्त मान्यताएँ रह सकती हैं। इसलिए अंतिम लक्ष्य के साथ-साथ अपनी मान्यताओं
का भी लेखा रखना आवश्यक है।

आइए देखें कि !!a{derivation} कैसे दी जा सकती है—इस !!{formula} की:
$\lexists[x][!C(x,b)]$, मान्यताओं $\lexists[x][(!A(x) 
\land !B(x))]$ और $\lforall[x][(!B(x) \lif !C(x,b))]$ से। सामान्य रीति से
आरम्भ करते हुए निष्कर्ष को सबसे नीचे लिखते हैं।
\begin{prooftree}
\AxiomC{}
\UnaryInfC{$\lexists[x][!C(x,b)]$}
\end{prooftree}

हमारे पास काम करने के लिए दो पूर्वाधार हैं। पहले का उपयोग करने के लिए—अर्थात्
$\lexists[x][!C(x, b)]$ की !!a{derivation}, $\lexists[x][(!A(x)
  \land !B(x))]$ से खोजने के लिए—\Elim{\lexists} नियम लगाएँगे। उस पर
अभिलाक्षणिक-चर शर्त लागू होती है, इसलिए यही नियम पहले लगाएँगे। इससे मिलता है:
\begin{prooftree}
\AxiomC{$\lexists[x][(!A(x) \land !B(x))]$}
\AxiomC{$\Discharge{!A(a) \land !B(a)}{1}$}
\DeduceC{$\lexists[x][!C(x,b)]$}
\DischargeRule{\Elim{\lexists}}{1}
\BinaryInfC{$\lexists[x][!C(x,b)]$}
\end{prooftree}
जिन दो मान्यताओं के साथ हम काम कर रहे हैं, उनमें~$!B$ उभयनिष्ठ है। इस बिंदु
पर \Elim{\land} लगाकर~$!B(a)$ को अलग निकालना उपयोगी हो सकता है।
\begin{prooftree}
\AxiomC{$\lexists[x][(!A(x) \land !B(x)])$}
\AxiomC{$\Discharge{!A(a) 
\land !B(a)}{1}$}
\RightLabel{\Elim{\land}}
\UnaryInfC{$!B(a)$}
\DeduceC{$\lexists[x][!C(x,b)]$}
\DischargeRule{\Elim{\lexists}}{1}
\BinaryInfC{$\lexists[x][!C(x,b)]$}
\end{prooftree}

काम के लिए दूसरी मान्यता है~$\lforall[x][(!B(x) \lif
  !C(x,b))]$। यहाँ अभिलाक्षणिक-चर शर्त नहीं है, इसलिए \Elim{\lforall} लगाकर
$x$ को !!{constant}~$a$ से उदाहरणित कर सकते हैं और
$!B(a) \lif !C(a, b)$ पा सकते हैं। अब हमारे पास $!B(a) \lif !C(a,b)$ और
$!B(a)$, दोनों हैं। अगला चरण \Elim{\lif} नियम का सीधा प्रयोग होना चाहिए।
\begin{prooftree}
\AxiomC{$\lexists[x][(!A(x) \land !B(x))]$}
\AxiomC{$\lforall[x][(!B(x) \lif !C(x,b))]$}
\RightLabel{\Elim{\lforall}}
\UnaryInfC{$!B(a) \lif !C(a,b)$}
\AxiomC{$\Discharge{!A(a) 
\land !B(a)}{1}$}
\RightLabel{\Elim{\land}}
\UnaryInfC{$!B(a)$}
\RightLabel{\Elim{\lif}}
\BinaryInfC{$!C(a,b)$}
\DeduceC{$\lexists[x][!C(x,b)]$}
\DischargeRule{\Elim{\lexists}}{1}
\insertBetweenHyps{\hspace{-5em}}
\BinaryInfC{$\lexists[x][!C(x,b)]$}
\end{prooftree}
हम लक्ष्य के बहुत निकट हैं!{} \Intro{\lexists} का एक प्रयोग और लक्ष्य मिल
जाता है।
\begin{prooftree}
\AxiomC{$\lexists[x][(!A(x) \land !B(x))]$}
\AxiomC{$\lforall[x][(!B(x) \lif !C(x,b))]$}
\RightLabel{\Elim{\lforall}}
\UnaryInfC{$!B(a) \lif !C(a,b)$}
\AxiomC{$\Discharge{!A(a) 
\land !B(a)}{1}$}
\RightLabel{\Elim{\land}}
\UnaryInfC{$!B(a)$}
\RightLabel{\Elim{\lif}}
\BinaryInfC{$!C(a,b)$}
\RightLabel{\Intro{\lexists}}
\UnaryInfC{$\lexists[x][!C(x,b)]$}
\DischargeRule{\Elim{\lexists}}{1}
\insertBetweenHyps{\hspace{-5em}}
\BinaryInfC{$\lexists[x][!C(x,b)]$}
\end{prooftree}
चूँकि हर चरण पर हमने सुनिश्चित किया कि अभिलाक्षणिक-चर शर्तों का उल्लंघन न हो,
इसलिए विश्वासपूर्वक कह सकते हैं कि यह सही !!{derivation} है।
\end{ex}

\begin{ex}
!!a{derivation} दीजिए—इस !!{formula} की:
$\lnot\lforall[x][!A(x)]$, मान्यताओं $\lforall[x][!A(x)] 
\lif \lexists[y][!B(y)]$ और $\lnot\lexists[y][!B(y)]$ से। सामान्य रीति से
आरम्भ करते हुए लक्ष्य !!{formula} को सबसे नीचे लिखते हैं।
\begin{prooftree}
\AxiomC{}
\UnaryInfC{$\lnot\lforall[x][!A(x)]$}
\end{prooftree}
!!{derivation} की अंतिम पंक्ति निषेध है, इसलिए \Intro{\lnot} लगाने का प्रयत्न
करते हैं। इसके लिए यह पता लगाना होगा कि अंतर्विरोध को कैसे !!{derive} करें।
\begin{prooftree}
\AxiomC{$\Discharge{\lforall[x][!A(x)]}{1}$}
\DeduceC{$\lfalse$}
\DischargeRule{\Intro{\lnot}}{1}
\UnaryInfC{$\lnot\lforall[x][!A(x)]$}
\end{prooftree}
यहाँ तक सब ठीक है। \Elim{\lforall} लगा सकते हैं, पर यह स्पष्ट नहीं कि उससे
लक्ष्य तक पहुँचने में सहायता मिलेगी। इसके बजाय अपनी एक मान्यता का उपयोग करें।
$\lforall[x][!A(x)] \lif \lexists[y][!B(y)]$ को
$\lforall[x][!A(x)]$ के साथ लेने पर \Elim{\lif} नियम लगा सकेंगे।
\begin{prooftree}
\AxiomC{$\lforall[x][!A(x)] \lif \lexists[y][!B(y)]$}
\AxiomC{$\Discharge{\lforall[x][!A(x)]}{1}$}
\RightLabel{\Elim{\lif}}
\BinaryInfC{$\lexists[y][!B(y)]$}
\DeduceC{$\lfalse$}
\DischargeRule{\Intro{\lnot}}{1}
\UnaryInfC{$\lnot\lforall[x][!A(x)]$}
\end{prooftree}
अब काम के लिए एक अंतिम मान्यता बची है। लगता है कि उसकी सहायता से
\Elim{\lnot} लगाकर अंतर्विरोध तक पहुँच सकेंगे।
\begin{prooftree}
\AxiomC{$\lnot\lexists[y][!B(y)]$}
\AxiomC{$\lforall[x][!A(x)] \lif \lexists[y][!B(y)]$}
\AxiomC{$\Discharge{\lforall[x][!A(x)]}{1}$}
\RightLabel{\Elim{\lif}}
\BinaryInfC{$\lexists[y][!B(y)]$}
\RightLabel{\Elim{\lnot}}
\BinaryInfC{$\lfalse$}
\DischargeRule{\Intro{\lnot}}{1}
\UnaryInfC{$\lnot\lforall[x][!A(x)]$}
\end{prooftree}
\end{ex}

\begin{prob}
निम्नलिखित को दिखाने वाली !!{derivation}s दीजिए:
\begin{enumerate}
\item $\Proves (\lforall[x][!A(x)] \land \lforall[y][!B(y)]) \lif
\lforall[z][(!A(z) \land !B(z))]$.
\item $\Proves (\lexists[x][!A(x)] \lor \lexists[y][!B(y)]) \lif
\lexists[z][(!A(z) \lor !B(z))]$.
\item $\lforall[x][(!A(x) \lif !B)] \Proves \lexists[y][!A(y)] \lif !B$.
\item $\lforall[x][\lnot !A(x)] \Proves \lnot\lexists[x][!A(x)]$.
\item $\Proves \lnot\lexists[x][!A(x)] \lif \lforall[x][\lnot !A(x)]$.
\item $\Proves \lnot\lexists[x][\lforall[y][((!A(x,y) \lif \lnot
!A(y,y)) \land (\lnot !A(y,y) \lif !A(x,y)))]]$.
\end{enumerate}
\end{prob}

\begin{prob}
निम्नलिखित को दिखाने वाली !!{derivation}s दीजिए:
\begin{enumerate}
\item $\Proves \lnot\lforall[x][!A(x)] \lif \lexists[x][\lnot!A(x)]$.
\item $(\lforall[x][!A(x)] \lif !B) \Proves \lexists[y][(!A(y) \lif !B)]$.
\item $\Proves \lexists[x][(!A(x) \lif \lforall[y][!A(y)])]$.
\end{enumerate}
(इन सभी के लिए $\FalseCl$~नियम आवश्यक है।)
\end{prob}

\end{document}
